\chapter{Prototyping and iteration}
\label{vol11:ch07}

\epigraph{The prototype's job is not to be the product; it is to teach the design team what the product must become.}{}

\chapmeta{Half-life: medium-life. Archetypes: optimisation, control. Exercise target: 20. Project track: hybrid.}

\noindent
Prototyping is the engineering activity of building a partial or simplified instance of a design to learn from it. The prototype is not the product; it is the design's intermediate physical manifestation, built specifically to surface what cannot be learned from analysis or simulation alone. The discipline of prototyping is the discipline of building the right prototype, learning the right things from it, and committing or revising based on what was learned. This chapter develops the discipline and the failure modes that follow from prototypes that became products without the substantive engineering between them.

\section{Why prototype}\pathtag{core}

Designs that proceed from concept to production without prototyping make commitments based on analysis, simulation, and judgment that turn out, with measurable frequency, to be inadequate. The inadequacies surface only when the design is in production or in service; the cost of correction at that stage is the high-stage-of-development cost that Chapter 1 discussed.

Prototyping surfaces inadequacies earlier:

\textbf{Things analyses cannot predict.} The analysis assumes ideal geometry; the actual part has manufacturing variation. The analysis assumes specified materials; the actual material has lot-to-lot variation. The analysis assumes idealised loading; the actual loading has spectrum content the analysis did not capture. The analysis is internally consistent within its assumptions; the prototype exposes the assumptions.

\textbf{Things simulations cannot predict.} The simulation embodies a model; the model has scope; behaviour outside the model's scope is not captured. The prototype operates in the full physical environment; behaviours that the model omitted surface as observations on the prototype.

\textbf{Things human users surface.} The user interacts with the prototype; the interaction reveals usability problems, mismatches between the design and the user's mental model, workarounds the design forces, errors the design enables. None of these surface in analysis or simulation alone.

\textbf{Things integration surfaces.} The prototype integrates the design's subsystems; the integration reveals interface mismatches, timing problems, compatibility issues, cumulative tolerance accumulation that the design's component-level analyses did not address.

\textbf{Things the test conditions surface.} The prototype is operated under conditions the design must handle in service; the test conditions surface failures the analyses had not predicted. The aviation industry's flight-test programs surface aerodynamic, structural, and control-system behaviours that the wind tunnel and the simulation did not capture; the post-flight-test design changes are routine.

\textbf{Things the team learns by building.} The act of building the prototype teaches the design team about the design's manufacturing characteristics, its assembly sequence, its maintainability, its packaging constraints. The learning informs subsequent design decisions and the engineer's growing intuition about the design class.

The economics of prototyping are favourable in most cases. The cost of a prototype is substantial (materials; engineering time; testing) but is reliably lower than the cost of discovering the prototype's lessons after production. The number of prototype iterations matters: a single prototype that confirms the design is adequate provides little information; multiple prototypes with deliberate variation provide the most learning per iteration.

\begin{principle}[Prototype to learn]
The prototype's purpose is to learn what the design's other activities cannot teach. Designs that proceed without prototyping commit to assumptions whose validity is unconfirmed; the assumptions are confirmed in service if they were correct or revealed as inadequate if they were not. The discipline of prototyping is the discipline of getting the confirmation or refutation at a stage when the design can still be revised cheaply.
\end{principle}

\section{Prototype types: foam, functional, beta, MVP}\pathtag{core}

Prototypes occupy a spectrum from very low fidelity (rough sketches in physical form) to nearly-production (manufacturing-representative units). The choice of fidelity depends on what the prototype is intended to teach.

\textbf{Foam models and low-fidelity mockups.} Physical representations of the design's form, built quickly from foam, cardboard, or 3D-printed material. The model has the design's outer shape but no internal function. The lessons it teaches: ergonomic fit; visual perception; size relative to the user and the environment; user reaction to form. The Apple industrial-design tradition uses foam mockups extensively at early conceptual stages; they are cheap to produce and fast to iterate on.

\textbf{Looks-like prototypes.} Physical representations that resemble the production design in form and surface finish but may have simplified internal structure. The lessons: aesthetic evaluation; marketing visualisation; user reaction to the finished appearance. The cost is higher than foam mockups; the value is the higher-fidelity user feedback they support.

\textbf{Works-like prototypes (functional).} Working implementations of the design's principal function, often with simplified packaging or non-final components. A circuit on a breadboard; a mechanism in clear plastic for visibility; a software prototype with placeholder UI. The lessons: function verification; performance measurement; integration of subsystems; failure mode identification.

\textbf{Looks-like-and-works-like prototypes.} Functional implementations in production-representative form. The lessons: integrated user evaluation; performance under realistic operating conditions; lifecycle behaviour over short test periods. The cost is high; the value is the most substantive pre-production design feedback.

\textbf{Engineering development units (EDUs).} Pre-production units built for engineering test purposes: environmental qualification; durability testing; certification testing. EDUs typically use production tooling but in limited quantities; the lessons are about the production design under controlled test conditions.

\textbf{Beta units.} Production-representative units released to selected users (often customers or partners) for use under realistic conditions. The lessons: usability over extended time; reliability under actual service; user satisfaction; issues that limited testing did not surface.

\textbf{Minimum Viable Product (MVP).} A version of the product with sufficient functionality to be released to early users; the MVP is intentionally minimal to enable rapid learning from real-world use before further investment. The MVP framing is dominant in software; it has analogues in hardware (limited initial production runs; pilot manufacturing) but the iteration speed is slower.

\textbf{Pilot production.} A limited production run using production tooling and processes; the run validates the production process, supports operator training, and produces units for further test or for limited market introduction.

The discipline of prototype-type selection is the discipline of matching the prototype to the learning the design needs. Foam mockups for ergonomic decisions; breadboard for circuit function; full-fidelity prototypes for integrated evaluation; pilot production for manufacturing readiness. Mismatching produces effort without learning: a foam mockup will not validate the circuit's function; a breadboard will not provide ergonomic feedback.

\section{Prototype-driven learning}\pathtag{core}

The prototype's value is the learning it produces. The learning happens through structured engagement with the prototype: not as an end in itself, but as a means to surface the design's behaviours and to inform the design's revision.

The activities that produce prototype-driven learning:

\textbf{Defined test plan.} Before the prototype is built, the design team defines what they intend to learn from it. The plan specifies: the questions the prototype will answer; the tests that will produce the answers; the measurements that will be taken; the acceptance criteria that will inform the design decisions. Without the plan, the prototype's testing is unstructured; the lessons are diffuse and the design's revisions are intuitive rather than evidence-based.

\textbf{Instrumented prototype.} The prototype is built with instrumentation that captures the data the test plan requires: sensors for force, displacement, temperature, voltage, current, vibration; logging for software behaviour; video for user interactions. The instrumentation is part of the prototype design, not an afterthought.

\textbf{Controlled test conditions.} The prototype is tested under conditions specified in the test plan: the inputs are controlled; the environment is controlled; the test sequence is repeatable. The control makes the observations interpretable; uncontrolled testing produces observations that conflate the prototype's behaviour with the test conditions' variability.

\textbf{User-engagement protocol.} If the prototype's testing involves users, the engagement is structured: the participants are recruited to represent the user population; the tasks are specified; the observations are documented; the user's verbal feedback is captured. The structure prevents the user testing from becoming a demonstration of the design's strengths and avoiding its weaknesses.

\textbf{Documented findings.} The prototype's behaviour and the testing's findings are documented in a form the design team can use for design revisions and that becomes part of the design's record. The documentation includes the unexpected findings as well as the expected ones; the unexpected are usually the most valuable.

\textbf{Design revisions.} The findings drive specific design revisions; the revisions are tracked back to the prototype's findings; the next prototype's test plan addresses what the previous prototype could not. The iteration is a loop with explicit purpose at each cycle.

The discipline addresses a common pattern: prototypes are built and operated, but the learning is not extracted systematically; the team's intuition is updated diffusely; the design's revisions are not traceable to specific findings. The discipline is the discipline of converting prototype experience into recorded, actionable engineering learning.

\section{Iteration cadence}\pathtag{standard}

The number and pacing of prototype iterations is part of the design's strategy. The classical pattern in product development is multiple iterations through escalating fidelity: low-fidelity early; medium-fidelity middle; high-fidelity late. The pattern's logic is that early iterations are cheap and answer many questions per dollar; late iterations are expensive and answer the questions that require high fidelity.

The cadence considerations:

\textbf{Iteration cycle time.} The time from one prototype to the next is the iteration cycle time. Software iterations can be daily; mechanical hardware iterations are typically weeks to months; full vehicle iterations are months to years. The cycle time bounds how many iterations the project can complete; faster cycle times enable more iterations and more learning per project.

\textbf{Iteration overlap.} Multiple prototypes can be in flight at once: one being built; another being tested; a third being analysed for findings. The overlap increases the effective throughput at the cost of complexity in managing the multiple work streams.

\textbf{Iteration depth.} The depth of design revision between iterations varies: small parameter changes; component substitutions; architectural changes. Deeper revisions usually require longer cycle times; the team's strategy is the balance between iteration speed and revision depth.

\textbf{Cost trajectory.} Early prototypes are cheap (low-fidelity; small lots); late prototypes are expensive (production-representative; tooling investment). The cumulative cost across iterations is substantial; the discipline includes deciding which iterations are necessary versus which would be incremental at high cost.

\textbf{Decision points.} Specific iterations are tied to project decisions: the iteration that supports the concept-selection decision; the iteration that supports the design-freeze decision; the iteration that supports the production-readiness decision. The iteration's findings must be available before the decision is required.

\textbf{Learning saturation.} At some point, additional iterations produce diminishing returns: the major lessons have been learned; subsequent iterations confirm rather than discover. The discipline includes recognising saturation and moving to the next stage rather than continuing iteration past the point of value.

\textbf{Schedule pressure.} Iteration count is often cut under schedule pressure. The cut is rational if the cut iterations would have produced diminishing returns; it is harmful if the cut iterations would have surfaced consequential issues. The discipline includes resisting the cut when the cut iterations are substantive and accepting it when they are not.

The software industry's continuous-integration / continuous-deployment (CI/CD) model is a high-cadence extreme: iterations may be per-commit (multiple per day); the prototype is the entire system; the testing is automated; the deployment is selective (canary releases to limited users) before broad release. The model is appropriate for software where the deployment cost is low and the rollback cost is low; it is partly transferable to hardware (e.g., over-the-air software updates for vehicles) but not fully (the physical hardware itself cannot be re-deployed at software cadence).

\section{Knowing when to commit}\pathtag{standard}

Iteration cannot continue indefinitely; at some point the design is committed to the next stage. The commitment is a judgment under uncertainty: the design's remaining uncertainty must be acceptable for the stage being committed to.

The commitment criteria depend on the stage:

\textbf{Concept commitment.} The design moves from concept generation to detailed design. The criteria: the concept is well-defined; the major requirements are addressed; the dominant risks are identified and have mitigation plans; the alternative concepts have been compared. The commitment is reversible (returning to concept generation is possible if detailed design reveals an issue) but expensive.

\textbf{Design freeze.} The detailed design is committed; subsequent changes are governed by formal change control. The criteria: the design meets the requirements (verified through analysis, simulation, and prototype); the manufacturing process is identified; the supply chain is established; the certification path is clear. The commitment is much harder to reverse; design freeze unlocks the manufacturing-tooling investment.

\textbf{Production readiness.} The design is committed to production. The criteria: the design has been validated against the requirements (Chapter 8); the manufacturing process has been qualified through pilot production; the supply chain is operational; the regulatory approvals are in hand. The commitment is the transition from design to manufacture; the design's properties are now determined by the as-built units rather than by ongoing design work.

\textbf{Release to service.} The product is delivered to users. The criteria: the production-readiness criteria plus customer acceptance, service-support infrastructure, training materials, and the operational decision to begin distribution. The commitment exposes the design to the full population of users and use conditions.

The discipline of commitment includes:

\textbf{Explicit criteria.} The commitment criteria are documented; the evidence against them is gathered; the commitment decision is made against the evidence. Tacit commitment ("we feel ready") produces inconsistent decisions; explicit commitment produces decisions that can be audited and improved.

\textbf{Stakeholder alignment.} The commitment requires sign-off from the parties whose subsequent work depends on it: manufacturing for design freeze; service for release; finance for budget release; safety for hazard acceptance; certification for regulatory path. The alignment surfaces objections that should be addressed before commitment.

\textbf{Risk acceptance.} The commitment includes explicit acceptance of the residual risks: known issues that will be addressed post-release; uncertainties that have not been fully resolved; reliability characteristics that will be confirmed in service. The acceptance is documented; the responsible decision-makers are recorded.

\textbf{Reversibility plan.} For commitments that may need to be reversed (recalls; service bulletins; design changes), the reversibility is planned in advance: how the recall would be managed; what the design-change pathway looks like; what the cost and timing are. The plan does not reduce the cost of reversal but it does reduce the surprise.

The Volume X failures often include premature commitment: the design was committed before the residual risks were acceptable for the commitment stage. The Boeing 737 MAX MCAS commitment to production with single-AoA architecture is an instance; the system safety assessment that should have warranted the commitment had not been rerun after the authority expansion. The discipline of commitment is the discipline of resisting these patterns.

\section{Pilot production and pre-production runs}\pathtag{standard}

Between full prototypes and full production lies the pilot run: a limited production using the production tooling and processes. The pilot's purpose is to validate the manufacturing process, surface manufacturing-related issues, train operators, and produce units for further test or for limited market introduction.

The pilot's activities:

\textbf{Tooling validation.} The production tooling (moulds; jigs; fixtures; assembly stations; test equipment) is exercised; tooling problems (dimensional accuracy; wear; cycle time) are surfaced and corrected.

\textbf{Process validation.} The production process is exercised end-to-end; process parameters (temperatures; pressures; times; sequence) are confirmed and adjusted; the as-produced units' conformance to specification is measured.

\textbf{Operator training.} The production operators are trained on the actual process; their experience reveals procedure clarifications, fixture improvements, and process refinements that the development engineers had not anticipated.

\textbf{Yield characterisation.} The pilot's yield (the fraction of units that pass all production checks) is measured; the yield's relationship to process parameters is characterised; the production process is adjusted to meet the yield target.

\textbf{First-article inspection.} The first units off the production tooling are inspected against the design specification in detail; the inspection confirms that the production process produces the designed part. Deviations are tracked back to tooling, process, or specification issues and addressed.

\textbf{Supply-chain qualification.} The supply chain (component suppliers; raw material suppliers; service providers) is exercised at production scale; capacity issues, quality issues, lead-time issues that pilot quantities surface are addressed before full production exposure.

\textbf{Documentation finalisation.} The manufacturing documentation (work instructions; quality plans; inspection procedures) is updated based on pilot experience; the documentation that will guide full production is finalised.

\textbf{Pre-production units.} The pilot's output units are used for: final certification testing; pre-launch customer evaluation; service-organisation training; spare-parts inventory for early service support. The units are the bridge between development and full production.

The pilot's discipline is to surface and address production-related issues before full production exposes them at much larger scale. A pilot that runs smoothly is sometimes a warning sign: it may indicate that the pilot's conditions were not representative of full production (e.g., the most experienced operators were assigned to the pilot; the most carefully-prepared materials were used). The pilot's value is in finding the issues; a pilot that finds nothing should be re-examined for whether it actually exercised the production at scale.

\section{Failure: prototypes that became products without the work in between}\pathtag{core}

The failure mode this section addresses is the prototype that was treated as the product: the design team built a prototype, it worked, and the prototype was put into production with minimal subsequent engineering. The pattern occurs in startups under capital pressure, in research-translation contexts, in proof-of-concept projects that gain commercial traction, and in time-pressured product launches.

The mechanisms by which prototype-to-product transitions fail:

\textbf{Pilot-scale assumptions in production.} The prototype was built one unit at a time with engineering attention; production builds many units with operator attention. The processes that work at pilot scale (manual fitting; selective assembly; engineering judgment at each step) do not work at production scale; production reveals consistency problems the prototype did not show.

\textbf{Tolerance accumulation.} The prototype was assembled with parts selected for fit; production produces parts with full tolerance range; the assembly's behaviour with full-tolerance parts differs from the prototype's behaviour with selected parts. The Hyatt Regency walkway case has elements of this: the design that was analysed differed from the design that was built; the difference was consequential.

\textbf{Reliability and durability untested.} The prototype was tested for function; the production design must operate reliably over its service life; the reliability and durability are different engineering questions that the prototype testing did not address. The Boeing 787 lithium-ion battery incidents (Volume X Chapter 4 case; here as Volume XI prototype-to-production reference) had elements: the prototype-stage thermal-runaway behaviour was characterised under conditions that did not fully represent service.

\textbf{Manufacturing variability not characterised.} The prototype was built by skilled engineers; production is by operators; the production variability was not characterised at prototype stage. Issues that the engineers handled implicitly (de-burring; alignment; lubrication) are not part of the production process; the production units have issues the prototype did not.

\textbf{User variability not characterised.} The prototype was tested by users similar to the design team; production reaches users with much broader variation; the design's behaviour with the broader user population was not characterised. Therac-25 has this characteristic: the operator-interface behaviour with the testing-stage operators differed from the behaviour with field operators under workflow pressure.

\textbf{Environmental envelope not characterised.} The prototype was tested in laboratory conditions; production units operate in the field's full environmental envelope (temperature; humidity; vibration; EMI; contamination); behaviours that the lab did not test surface in service.

\textbf{Service-life behaviour not characterised.} The prototype was tested for hours or days; production units operate for years; degradation, fatigue, wear, and aging are not characterised. The Aloha Airlines 243 fuselage fatigue (Volume X Chapter 2) has elements: the long-term cumulative fatigue and disbond behaviour was not adequately characterised before the design proceeded.

\textbf{Regulatory and certification gaps.} The prototype was not subjected to full regulatory testing; production exposes the design to the full regulatory framework; certification issues surface that the prototype testing did not. The Volkswagen diesel emissions case has elements: the certification testing produced compliant results; the production units' actual emissions did not.

The remediation is the substantive engineering work between prototype and production: the pilot run, the qualification testing, the certification testing, the reliability characterisation, the manufacturing-process validation, the user validation at the broader population, the service-life testing. The work takes time and resources; the alternative (skipping the work) is the failure mode this section addresses.

The discipline of project management includes the discipline of resisting the pressure to compress prototype-to-product activities below what the design's actual risks require. The compression is reliably cheaper in the moment and reliably more expensive in the aggregate, when the issues that the compressed work would have caught surface in service.

\begin{archetype}[control]
Prototyping is the control archetype applied to design uncertainty. The prototype is the design's measurement against reality; the iteration is the control loop; the design's commitment criteria are the stability requirements. The discipline of iteration is the discipline of staying in the control loop long enough to reduce the design's uncertainty to acceptable levels before commitment.
\end{archetype}

\begin{principle}[The prototype is not the product]
The prototype's purpose is learning; the product's purpose is service. The engineering work between prototype and product is the work that transforms a learning instance into a serviceable artefact: pilot production, qualification, validation, manufacturing-process development, service-life characterisation. The work is necessary; designs that skip it produce products whose failures had been latent in the prototype-to-product gap.
\end{principle}

\begin{project}[Three prototype iterations]
Select a small device to design (a mechanism; a circuit; a software tool; a household helper). Develop and build three prototype iterations, with a defined learning objective for each iteration and a documented update from each to the next. Iteration 1: rough proof-of-concept; demonstrate principal function. Iteration 2: address two specific issues identified in Iteration 1 and add one element of usability evaluation. Iteration 3: produce a near-final version, with documented test results against requirements you defined at the start. Deliverable: brief documentation of each iteration (one page each), the test results, and a one-page reflection on what the iterations taught you that analysis or simulation would not have.

\textbf{Hazard.} The temptation will be to make Iteration 1 already pretty good and to make Iterations 2 and 3 incremental refinements. The discipline is the opposite: make Iteration 1 deliberately rough so that Iterations 2 and 3 have something substantive to learn from. The point is the iteration discipline, not the polish of the final artefact.
\end{project}

\section{Exercises}

\subsection*{Calculation}\pathtag{core}

\begin{exercise}
A prototype iteration takes 2 weeks (build) plus 1 week (test) plus 1 week (design revision). The development budget allows 6 months between concept commitment and design freeze. Compute the maximum number of iterations the budget supports. If parallel iterations (one being built while another is being tested) are possible, compute the increased iteration count.
\end{exercise}

\begin{exercise}
A prototype's first-pass yield is 60\%; the second-iteration yield is 80\%; the third-iteration yield is 95\%; full production must achieve $\geq 99\%$. Compute the cost per acceptable unit at each stage if the per-unit fabrication cost is constant. Discuss the economic argument for pilot-stage iteration.
\end{exercise}

\begin{exercise}
A prototype test programme identifies 12 issues; addressing each costs (on average) 40 engineering hours; the rate of new issue discovery decreases by half per iteration. Compute the cumulative engineering hours through 5 iterations.
\end{exercise}

\subsection*{Derivation}\pathtag{standard}

\begin{exercise}
Derive a simple model of the relationship between prototype iteration count and post-release defect rate. Assume each iteration catches a fraction $\alpha$ of remaining defects; derive the post-iteration-$N$ defect count as a function of the initial defect count and $\alpha$. Identify the iteration count at which the marginal benefit of an additional iteration falls below a specified threshold.
\end{exercise}

\begin{exercise}
For a multi-fidelity prototyping strategy (low-fidelity for early iterations; high-fidelity for late iterations), derive the cost-optimal mix of iterations as a function of the cost ratio between fidelity levels and the issue-discovery probability per fidelity level.
\end{exercise}

\subsection*{Estimation}\pathtag{standard}

\begin{exercise}
Estimate the cost of a complete prototype-to-production engineering programme for a mid-complexity product (e.g., a consumer electronic device): low-fidelity prototypes; functional prototypes; EDUs; pilot production. Identify the fraction of the cost at each stage.
\end{exercise}

\begin{exercise}
Estimate the cost of a major recall (e.g., a vehicle recall affecting 100,000 units) and compare to the cost of the prototype-to-production engineering that would have caught the underlying issue.
\end{exercise}

\begin{exercise}
Estimate the typical iteration count in a software product (where iterations may be measured in commits and continuous deployment is the norm) versus a vehicle product (where iterations may be measured in build-test-revise cycles).
\end{exercise}

\subsection*{Simulation}\pathtag{standard}

\begin{exercise}
Implement a simulation of the iteration-saturation phenomenon: each iteration discovers issues at rate proportional to the remaining unknown issues; the total issue count is bounded. Sweep parameters and identify the iteration regime in which marginal returns diminish below a threshold.
\end{exercise}

\begin{exercise}
Simulate a stochastic-yield process: the prototype's yield improves with iteration count according to a learning curve; each iteration costs a fixed amount; the production cost depends on the final yield. Find the iteration count that minimises total cost (iteration cost plus production cost).
\end{exercise}

\subsection*{Design}\pathtag{standard}

\begin{exercise}
Design a prototype-test plan template for an engineering project. Include: learning objectives; test conditions; instrumentation; success criteria; documentation requirements; design-revision triggers.
\end{exercise}

\begin{exercise}
Design a pilot-production plan for a new product. Specify: pilot quantity; tooling validation; process validation; operator training; first-article inspection; yield characterisation; transition criteria to full production.
\end{exercise}

\begin{exercise}
Design a commitment-criteria framework for an engineering organisation's design process. Specify the criteria for each stage commitment (concept commitment; design freeze; production readiness; release to service) and the documentation each requires.
\end{exercise}

\subsection*{Judgement}\pathtag{mastery}

\begin{exercise}
Argue: the MVP (Minimum Viable Product) framing is appropriate for software but misleading for hardware. Identify the structural differences between software and hardware iteration that bear on the dispute and discuss the implications for hardware-product development under MVP-style management.
\end{exercise}

\begin{exercise}
Discuss the prototype-to-production failure pattern in the context of startup engineering. Argue: capital pressure on startups produces structural pressure to skip prototype-to-production engineering. Argue the opposite: startup constraints can be addressed without skipping the substantive engineering. Identify the conditions that determine which is correct.
\end{exercise}

\begin{exercise}
The Boeing 787 lithium-ion battery incidents (Volume X reference) had prototype-stage thermal-runaway behaviour that was not adequately characterised before commitment to production. Discuss what prototype-engineering activities would have caught the inadequacy and the institutional factors that may have prevented them.
\end{exercise}

\begin{exercise}
A design team's prototype works in the laboratory under controlled conditions. The product manager asks "can we ship it?" Discuss the engineering response: what additional activities are required before commitment to production; how to communicate the requirement to the product manager; the institutional structures that support the engineering discipline against schedule pressure.
\end{exercise}

\begin{exercise}
Argue: pilot production that "works perfectly the first time" should be treated as a warning sign rather than as a success indicator. Discuss the structural reasons that pilots that find no issues may be pilots that did not actually exercise the production at scale, and identify the diagnostic questions that distinguish substantive from theatre pilots.
\end{exercise}

\begin{exercise}
A medical-device design has been prototyped successfully in the laboratory. The next stage requires evaluation by clinicians in actual clinical workflows. Discuss the design of the clinical evaluation: which clinicians (the most experienced; the average; the most-burdened); which clinical settings (academic medical centre; community hospital; emergency department); which patient populations; the criteria for the design to proceed to regulatory submission versus to return for further iteration.
\end{exercise}

\begin{exercise}
A startup has built a prototype that demonstrates the principal function of their product. Investors are pressing for product launch within 6 months; the engineering team estimates that the prototype-to-production engineering would take 18 months under normal practice. Discuss the engineering team's options for handling the pressure: which prototype-to-production activities are most essential; which can be deferred to early field operation; which require external (regulatory) sign-off that cannot be deferred. Identify the professional-ethics framework that bears on the situation.
\end{exercise}
